\chapter[16]{Perelman's Differential Harnack Estimate}
\label{ch:chow-ii-perelman-differential-harnack}
\dcref{ch16}
\source{chow-et-al-ricci-flow-part-ii:page-0332}

Use backward time \(\tau\) and write \(d\mu\) for the Riemannian measure.

\section{The heat equation}

\begin{definition}[heat-kernel logarithmic density]
\label{def:chow-ii-ch16-heat-log-density}
\source{chow-et-al-ricci-flow-part-ii:page-0333}
\dcref{ch16:eq-16.5}
\group{Heat Harnack}

For a positive heat solution $u$ on $(M^n,g)$ and $\tau>0$, set
\[
u=(4\pi\tau)^{-n/2}e^{-f}.
\]
Then $\partial_\tau f=\Delta f-|\nabla f|^2-n/(2\tau)$.
\end{definition}

\begin{theorem}[Li--Yau differential Harnack estimate]
\label{thm:chow-ii-ch16-li-yau}
\source{chow-et-al-ricci-flow-part-ii:page-0334}
\uses{def:chow-ii-ch16-heat-log-density}
\dcref{ch16:16.1}
\group{Heat Harnack}


If $\operatorname{Rc}\ge0$, then every positive heat solution satisfies
\[
\partial_\tau f+|\nabla f|^2-\Delta f+\frac{n}{2\tau}=0,
\qquad \Delta f-\frac{n}{2\tau}\le0.
\]
Consequently, for $\tau_2>\tau_1$,
\[
f(x_2,\tau_2)-f(x_1,\tau_1)\le
 \frac{d(x_1,x_2)^2}{4(\tau_2-\tau_1)}.
\]
\end{theorem}
\begin{proof}
\sketch The Bochner formula applied to the equation in
Definition~\ref{def:chow-ii-ch16-heat-log-density}, followed by the scalar
maximum principle, gives the differential inequality. Integrating it along a
space-time path and minimizing the weighted kinetic energy gives the last
display.
\end{proof}

\begin{theorem}[Cheeger--Yau lower bound]
\label{thm:chow-ii-ch16-cheeger-yau}
\source{chow-et-al-ricci-flow-part-ii:page-0334}
\dcref{ch16:16.3}
\group{Heat Kernel Bounds}

If $(M^n,g)$ is complete with $\operatorname{Rc}\ge0$, its positive heat kernel
satisfies
\[
H(x,y,\tau)\ge(4\pi\tau)^{-n/2}
 \exp\!\left(-\frac{d(x,y)^2}{4\tau}\right).
\]
\end{theorem}

\begin{theorem}[Li--Yau--Perelman estimate for the heat kernel]
\label{thm:chow-ii-ch16-heat-kernel-harnack}
\source{chow-et-al-ricci-flow-part-ii:page-0335}
\uses{def:chow-ii-ch16-heat-log-density, thm:chow-ii-ch16-cheeger-yau}
\dcref{ch16:16.4}
\group{Heat Harnack}


Let $H=(4\pi\tau)^{-n/2}e^{-f}$ be the positive fundamental solution on a
closed manifold with $\operatorname{Rc}\ge0$. Then
\[
\tau(2\Delta f-|\nabla f|^2)+f-n\le0.
\]
Equivalently, $2\tau\partial_\tau f+\tau|\nabla f|^2+f\le0$.
\end{theorem}
\begin{proof}
\sketch For $W=\tau(2\Delta f-|\nabla f|^2)+f-n$, the Bochner formula gives
\[
(\partial_\tau-\Delta)W=-2\langle\nabla f,\nabla W\rangle
 -2\tau\left(\left|\nabla^2f-\frac{g}{2\tau}\right|^2
 +\operatorname{Rc}(\nabla f,\nabla f)\right).
\]
The maximum principle and the short-time heat-kernel asymptotics imply
$W\le0$.
\end{proof}

\begin{definition}[linear entropy]
\label{def:chow-ii-ch16-linear-entropy}
\source{chow-et-al-ricci-flow-part-ii:page-0336}
\dcref{ch16:16.20}
\group{Linear Entropy}

For $\int_M(4\pi\tau)^{-n/2}e^{-f}\,d\mu=1$, define
\[
\mathcal W_{\rm lin}(g,f,\tau)=
 \int_M\bigl(\tau(2\Delta f-|\nabla f|^2)+f-n\bigr)
 \frac{e^{-f}}{(4\pi\tau)^{n/2}}\,d\mu.
\]
Integration by parts gives
\[
\mathcal W_{\rm lin}=
 \int_M(\tau|\nabla f|^2+f-n)\frac{e^{-f}}{(4\pi\tau)^{n/2}}\,d\mu.
\]
\end{definition}

\begin{theorem}[linear entropy formula]
\label{thm:chow-ii-ch16-linear-entropy-formula}
\source{chow-et-al-ricci-flow-part-ii:page-0337,chow-et-al-ricci-flow-part-ii:page-0340}
\uses{def:chow-ii-ch16-linear-entropy, def:chow-ii-ch16-heat-log-density}
\dcref{ch16:16.8}
\group{Linear Entropy}


Let $u=(4\pi\sigma)^{-n/2}e^{-f}$ solve $\partial_\tau u=\Delta u$ on a
closed manifold, with $\sigma'=1$ and $\int_Mu\,d\mu=1$. Then
\[
\frac{d}{d\tau}\mathcal W_{\rm lin}(f,\sigma)
 =-2\sigma\int_M\left(\left|\nabla^2f-\frac{g}{2\sigma}\right|^2
 +\operatorname{Rc}(\nabla f,\nabla f)\right)u\,d\mu.
\]
In particular, nonnegative Ricci curvature makes $\mathcal W_{\rm lin}$
nonincreasing.
\end{theorem}
\begin{proof}
\sketch Put $\omega=2\Delta f-|\nabla f|^2$ and $W=\sigma\omega+f-n$.
Bochner's identity gives
\[
(\partial_\tau-\Delta)\omega=-2\langle\nabla f,\nabla\omega\rangle
 -2(|\nabla^2f|^2+\operatorname{Rc}(\nabla f,\nabla f)),
\]
and completing the square gives the corresponding equation for $W$. Multiplying
by $u$, integrating by parts, and using $\partial_\tau u=\Delta u$ proves the
formula.
\end{proof}

\section{Complete manifolds and heat-kernel estimates}

\begin{theorem}[heat-kernel parametrix]
\label{thm:chow-ii-ch16-heat-parametrix}
\source{chow-et-al-ricci-flow-part-ii:page-0341}
\dcref{ch16:16.12}
\group{Heat Kernel Bounds}

For a complete manifold and its minimal positive heat kernel,
\[
H(x,y,\tau)\sim(4\pi\tau)^{-n/2}e^{-d(x,y)^2/(4\tau)}
 \sum_{j=0}^{\infty}\tau^ju_j(x,y,\tau),
\]
locally uniformly with the standard remainder estimates and
$u_0(x,x,0)=1$.
\end{theorem}

\begin{definition}[asymptotic volume ratio]
\label{def:chow-ii-ch16-avr}
\source{chow-et-al-ricci-flow-part-ii:page-0342}
\dcref{ch16:16.14}
\group{Heat Kernel Bounds}

For a complete noncompact manifold with $\operatorname{Rc}\ge0$, set
\[
\operatorname{AVR}(g)=\lim_{r\to\infty}
 \frac{\operatorname{Vol}B(p,r)}{\omega_nr^n}.
\]
\end{definition}

\begin{theorem}[Li--Yau upper heat-kernel bound]
\label{thm:chow-ii-ch16-heat-upper-bound}
\source{chow-et-al-ricci-flow-part-ii:page-0342}
\uses{def:chow-ii-ch16-avr}
\dcref{ch16:16.13}
\group{Heat Kernel Bounds}


If $\operatorname{Rc}\ge0$, then for every $\delta\in(0,1)$,
\[
H(x,y,\tau)\le\frac{C(n,\delta)}{\operatorname{Vol}B(x,\sqrt\tau)}
 \exp\!\left(-\frac{d(x,y)^2}{(4+\delta)\tau}\right).
\]
If $\operatorname{AVR}(g)>0$, the denominator may be replaced by
$\omega_n\operatorname{AVR}(g)\tau^{n/2}$.
\end{theorem}

\begin{corollary}[finite heat-kernel entropy and gradient bound]
\label{cor:chow-ii-ch16-heat-integrability}
\source{chow-et-al-ricci-flow-part-ii:page-0342,chow-et-al-ricci-flow-part-ii:page-0344}
\uses{thm:chow-ii-ch16-heat-upper-bound, thm:chow-ii-ch16-cheeger-yau}
\dcref{ch16:16.15}
\group{Heat Kernel Bounds}


For $H=(4\pi\tau)^{-n/2}e^{-f}$ on a complete manifold with
$\operatorname{Rc}\ge0$,
\[
\int_M fH\,d\mu\le C(n),\qquad
|\nabla f|(x,y,\tau)\le C(n)\tau^{-1/2}\left(1+\frac{d(x,y)^2}{\tau}\right).
\]
Consequently $\mathcal W_{\rm lin}(f,\tau)$ is finite and the integrations by
parts used in its formula are valid.
\end{corollary}

\begin{theorem}[Hamilton gradient and Laplacian estimates]
\label{thm:chow-ii-ch16-hamilton-heat-estimates}
\source{chow-et-al-ricci-flow-part-ii:page-0345,chow-et-al-ricci-flow-part-ii:page-0347}
\dcref{ch16:16.18}
\group{Heat Kernel Bounds}

Let $u>0$ solve the heat equation on a complete manifold with
$\operatorname{Rc}\ge0$ and $u\le A$. Then
\[
\tau\frac{|\nabla u|^2}{u^2}\le\log\frac Au,\qquad
\tau\left(\frac{\Delta u}{u}+\frac{|\nabla u|^2}{u^2}\right)
 \le n+4\log\frac Au.
\]
\end{theorem}

\section{Entropy rigidity and localization}

\begin{theorem}[heat-kernel differential Harnack on complete manifolds]
\label{thm:chow-ii-ch16-complete-heat-harnack}
\source{chow-et-al-ricci-flow-part-ii:page-0352,chow-et-al-ricci-flow-part-ii:page-0355}
\uses{thm:chow-ii-ch16-linear-entropy-formula, cor:chow-ii-ch16-heat-integrability}
\dcref{ch16:16.28}
\group{Heat Harnack}


If $M$ is complete with $\operatorname{Rc}\ge0$ and $H$ is the heat kernel,
then, writing $H=(4\pi\tau)^{-n/2}e^{-f}$,
\[
\tau(2\Delta f-|\nabla f|^2)+f-n\le0,
\qquad \mathcal W_{\rm lin}(f,\tau)\le0.
\]
\end{theorem}
\begin{proof}
\sketch For $v_H=(\tau(2\Delta f-|\nabla f|^2)+f-n)H$ one has
\[
(\partial_\tau-\Delta)v_H=-2\tau\left(\left|\nabla^2f-\frac{g}{2\tau}\right|^2
 +\operatorname{Rc}(\nabla f,\nabla f)\right)H\le0.
\]
Pairing with a nonnegative backward heat solution and using the parametrix to
show that the pairing tends to zero as $\tau\downarrow0$ gives $v_H\le0$.
\end{proof}

\begin{corollary}[entropy rigidity]
\label{cor:chow-ii-ch16-entropy-rigidity}
\source{chow-et-al-ricci-flow-part-ii:page-0357}
\uses{thm:chow-ii-ch16-complete-heat-harnack, thm:chow-ii-ch16-linear-entropy-formula}
\dcref{ch16:16.31}
\group{Entropy Rigidity}


If $\operatorname{Rc}\ge0$ and $\mathcal W_{\rm lin}(f,\tau_0)\ge0$ for the
heat kernel at some $\tau_0>0$, then $(M,g)$ is isometric to $\mathbb R^n$.
\end{corollary}
\begin{proof}
\sketch The two preceding results force $\nabla^2f=(2\tau)^{-1}g$ for
$0<\tau\le\tau_0$. Thus $\phi=4\tau f$ satisfies $\nabla^2\phi=2g$;
its minimum is a pole and $\phi-\min\phi=d^2$. Hence $\Delta d^2=2n$,
and equality in Bishop--Gromov comparison gives Euclidean space.
\end{proof}

\begin{definition}[Euclidean logarithmic Sobolev inequality]
\label{def:chow-ii-ch16-euclidean-log-sobolev}
\source{chow-et-al-ricci-flow-part-ii:page-0358}
\dcref{ch16:16.33}
\group{Entropy Rigidity}

The Euclidean logarithmic Sobolev inequality is
\[
\int_M\left(\frac12|\nabla f|^2+f-n\right)u\,d\mu\ge0
\]
for every $u=(2\pi)^{-n/2}e^{-f}$ with $\int_Mu\,d\mu=1$.
\end{definition}

\begin{theorem}[Bakry--Concordet--Ledoux rigidity]
\label{thm:chow-ii-ch16-bakry-concordet-ledoux}
\source{chow-et-al-ricci-flow-part-ii:page-0358}
\uses{def:chow-ii-ch16-euclidean-log-sobolev, cor:chow-ii-ch16-entropy-rigidity}
\dcref{ch16:16.35}
\group{Entropy Rigidity}


On a complete manifold with $\operatorname{Rc}\ge0$, the Euclidean logarithmic
Sobolev inequality holds for every normalized positive density if and only if
$(M,g)$ is isometric to $\mathbb R^n$.
\end{theorem}

\begin{theorem}[linear entropy and asymptotic volume ratio]
\label{thm:chow-ii-ch16-entropy-avr}
\source{chow-et-al-ricci-flow-part-ii:page-0359,chow-et-al-ricci-flow-part-ii:page-0360}
\uses{def:chow-ii-ch16-avr, def:chow-ii-ch16-linear-entropy}
\dcref{ch16:16.39}
\group{Entropy Rigidity}


For a complete noncompact manifold with $\operatorname{Rc}\ge0$ and maximum
volume growth, the heat-kernel entropy has a limit and
\[
\lim_{\tau\to\infty}\mathcal W_{\rm lin}(f(\tau),\tau)=\log\operatorname{AVR}(g).
\]
\end{theorem}

\begin{definition}[parabolic heat-kernel cutoff]
\label{def:chow-ii-ch16-heat-cutoff}
\source{chow-et-al-ricci-flow-part-ii:page-0361}
\dcref{ch16:16.40}
\group{Localized Entropy}

For $\operatorname{Rc}\ge0$ define
\[
\varphi_{(x_0,t_0),\rho}(y,t)=\left(1-
 \frac{d(x_0,y)^2+2n(t-t_0)}{\rho^2}\right)_+.
\]
Then $(\partial_t-\Delta)\varphi_{(x_0,t_0),\rho}\le0$ distributionally.
\end{definition}

\begin{theorem}[localized linear entropy formula]
\label{thm:chow-ii-ch16-local-linear-entropy}
\source{chow-et-al-ricci-flow-part-ii:page-0361,chow-et-al-ricci-flow-part-ii:page-0362}
\uses{def:chow-ii-ch16-heat-cutoff, thm:chow-ii-ch16-complete-heat-harnack}
\dcref{ch16:16.42}
\group{Localized Entropy}


With $v_H$ as above and $\tau=t_0-t$,
\[
\frac{d}{dt}\int_M\varphi_{(x_0,t_0),\rho}(-v_H)\,d\mu
\le-2\tau\int_M\varphi_{(x_0,t_0),\rho}
 \left(\left|\nabla^2f-\frac{g}{2\tau}\right|^2
 +\operatorname{Rc}(\nabla f,\nabla f)\right)H\,d\mu.
\]
\end{theorem}

\section{Perelman's estimate for Ricci flow}

\begin{definition}[adjoint heat fundamental solution]
\label{def:chow-ii-ch16-adjoint-heat-kernel}
\source{chow-et-al-ricci-flow-part-ii:page-0362}
\dcref{ch16:16.43}
\group{Perelman Harnack}

For a Ricci flow $g(t)$ and $\tau=T-t$, an adjoint heat kernel centered at
$(x,0)$ solves
\[
\left(\partial_\tau-\Delta+R\right)H=0,
\qquad H(\cdot,\tau)\longrightarrow\delta_x\quad(\tau\downarrow0).
\]
Writing $H=(4\pi\tau)^{-n/2}e^{-f}$, set
\[
v_H=\left(\tau(R+2\Delta f-|\nabla f|^2)+f-n\right)H.
\]
\end{definition}

\begin{theorem}[Perelman's differential Harnack estimate]
\label{thm:chow-ii-ch16-perelman-harnack}
\source{chow-et-al-ricci-flow-part-ii:page-0362,chow-et-al-ricci-flow-part-ii:page-0364}
\uses{def:chow-ii-ch16-adjoint-heat-kernel}
\dcref{ch16:16.44}
\group{Perelman Harnack}


For every Ricci flow on a closed manifold and its adjoint heat kernel,
\[
v_H\le0,\qquad
\tau(R+2\Delta f-|\nabla f|^2)+f-n\le0.
\]
Moreover, for any positive adjoint heat solution $u$,
\[
\left(\partial_\tau-\Delta+R\right)v_u
 =-2\tau\left|\operatorname{Rc}+\nabla^2f-\frac{g}{2\tau}\right|^2u.
\]
\end{theorem}
\begin{proof}
\sketch The displayed evolution identity follows by differentiating the adjoint heat
equation and commuting covariant derivatives. Pairing it with a nonnegative
solution of the conjugate forward heat equation, the adjoint heat parametrix
and the weak maximum principle imply that the pairing tends to zero at
$\tau=0$ and hence $v_H\le0$ pointwise.
\end{proof}

\begin{theorem}[adjoint heat-kernel parametrix]
\label{thm:chow-ii-ch16-adjoint-parametrix}
\source{chow-et-al-ricci-flow-part-ii:page-0364}
\uses{def:chow-ii-ch16-adjoint-heat-kernel}
\dcref{ch16:16.46}
\group{Perelman Harnack}


On a smooth Ricci flow, the adjoint heat kernel has the local expansion
\[
H(x,y,\tau)=(4\pi\tau)^{-n/2}e^{-d_{g(0)}(x,y)^2/(4\tau)}
 (u_0(x,y,0)+O(\tau)),\qquad u_0(x,x,0)=1.
\]
\end{theorem}

\begin{lemma}[reduced-distance subsolution]
\label{lem:chow-ii-ch16-reduced-distance-subsolution}
\source{chow-et-al-ricci-flow-part-ii:page-0370,chow-et-al-ricci-flow-part-ii:page-0371}
\dcref{ch16:16.48}
\group{Perelman Harnack}

For the reduced distance $\ell(y,\tau;x,0)$ under two-sided Ricci bounds,
\[
e^{-2k_1\tau}\frac{d_0(x,y)^2}{4\tau}-\frac{nk_1}{3}\tau
\le\ell\le e^{2k_2\tau}\frac{d_0(x,y)^2}{4\tau}+\frac{nk_2}{3}\tau,
\]
and
\[
\left(\partial_\tau-\Delta+R\right)
 \left((4\pi\tau)^{-n/2}e^{-\ell}\right)\le0
\]
weakly.
\end{lemma}

\begin{lemma}[adjoint kernel dominates reduced-distance kernel]
\label{lem:chow-ii-ch16-adjoint-dominates-reduced}
\source{chow-et-al-ricci-flow-part-ii:page-0372}
\uses{def:chow-ii-ch16-adjoint-heat-kernel, lem:chow-ii-ch16-reduced-distance-subsolution}
\dcref{ch16:16.49}
\group{Perelman Harnack}


For the fundamental solution centered at $(x,0)$,
\[
H(y,\tau;x,0)\ge(4\pi\tau)^{-n/2}e^{-\ell(y,\tau;x,0)},
\qquad f(y,\tau;x,0)\le\ell(y,\tau;x,0).
\]
\end{lemma}
\begin{proof}
\sketch The difference of the two kernels is an adjoint supersolution with zero
distributional initial data. Pairing it with arbitrary nonnegative solutions
of the conjugate forward heat equation and applying the weak maximum principle
proves nonnegativity.
\end{proof}

\begin{theorem}[Perelman reduced-distance Harnack estimate]
\label{thm:chow-ii-ch16-reduced-distance-harnack}
\source{chow-et-al-ricci-flow-part-ii:page-0379,chow-et-al-ricci-flow-part-ii:page-0380}
\uses{def:chow-ii-ch16-adjoint-heat-kernel, lem:chow-ii-ch16-adjoint-dominates-reduced}
\dcref{ch16:16.54}
\group{Perelman Harnack}


For a backward Ricci flow $\partial_\tau g=2\operatorname{Rc}$ on a closed
manifold, the adjoint heat kernel satisfies
\[
f(q,\tau)\le\ell(q,\tau):=\frac1{2\sqrt\tau}
 \inf_\gamma\int_0^\tau\sqrt{\bar\tau}
 \left(R+|\dot\gamma|^2\right)\,d\bar\tau.
\]
Equivalently, $H(q,\tau)\ge(4\pi\tau)^{-n/2}e^{-\ell(q,\tau)}$.
\end{theorem}

\begin{theorem}[localized Ricci-flow entropy formulas]
\label{thm:chow-ii-ch16-localized-perelman-entropy}
\source{chow-et-al-ricci-flow-part-ii:page-0380,chow-et-al-ricci-flow-part-ii:page-0381}
\uses{thm:chow-ii-ch16-perelman-harnack}
\dcref{ch16:16.56}
\group{Localized Entropy}


Let $\bar L=(t_1-t)\ell$ be the backward reduced distance and
\[
\psi=\left(1-\frac{\bar L+2n(t-t_2)}{\rho^2}\right)_+.
\]
Then
\[
\frac{d}{dt}\int_M(-v_H)\psi\,d\mu_t
\le-2\int_M\tau\left|\operatorname{Rc}+\nabla^2f-\frac{g}{2\tau}\right|^2H\psi\,d\mu_t.
\]
The same inequality holds with the forward reduced distance and its cutoff in
place of $\ell$ and $\psi$.
\end{theorem}
